\section{Related Work}
\label{sec:related_work}

Our work operates at the intersection of parameter-efficient fine-tuning, dynamic test-time model merging, stateful routing controllers, and multi-tenant deep learning serving infrastructures.

\paragraph{Parameter-Efficient Fine-Tuning and Model Merging.}
To avoid the prohibitive costs of training and deploying fully distinct copies of large-scale foundation models, Parameter-Efficient Fine-Tuning (PEFT) has become the standard. Low-Rank Adaptation (LoRA) \cite{hu2021lora} freezes the base model parameters and injects compact, trainable low-rank adapters. To combine capabilities from diverse task-specific adapters, model merging techniques seek to blend distinct parameters or representations. Weight-space merging methods, such as task arithmetic \cite{ilharco2022editing} and Git Re-Basin \cite{ainsworth2022git}, interpolate fine-tuned weights statically. While static merging provides a single unified model with zero latency overhead, it is unable to adapt to dynamically shifting workloads at inference time and suffers from representational collapse under severe task interference.

\paragraph{Stateless vs. Stateful Dynamic Routing.}
To address the limitations of static merging, test-time dynamic model merging has been proposed. Under these frameworks, ensembling coefficients $\boldsymbol{\alpha}_t$ are computed sample-by-sample. Sample-wise Activation Blending (SABLE) \cite{liang23sable} extracts intermediate activations at early boundary layers, projects them onto task-specific subspaces to form task coordinates, and applies a direct softmax mapping to compute ensembling weights. While stateless routing is highly responsive, it is highly sensitive to sample-level noise and outliers. This leads to a severe \emph{routing jitter paradox}, where high-frequency switching between experts across layers destabilizes the representation flow, degrading classification performance.

To suppress this jitter, stateful ensembling treatment has recently emerged. ChemMerge \cite{chemmerge25} models representation blending as continuous-time chemical kinetics reaction ODEs, acting as a temporal low-pass filter over raw activation coordinate signals. Similarly, PAC-Kinetics \cite{pac_zca_2026} formalizes this stateful smoothing as a stable first-order diagonal recurrence, utilizing a Catoni-type PAC-Bayesian bound to optimize kinetics and routing parameters on a calibration set. While these stateful routers dramatically reduce routing jitter, they rely heavily on the assumption that query streams are homogeneous. When query streams are interleaved (as in realistic deployment scenarios), their rigid global temporal smoothing leads to state contamination and catastrophic lag, a limitation we address in this paper.

\paragraph{Multi-Tenant Serving Architectures.}
In modern cloud infrastructures, deep learning models are served under massive multi-tenant workloads. Systems like Punica \cite{chen2023punica}, S-LoRA \cite{s-lora}, and vLLM \cite{kwon2023efficient} manage the scheduling and execution of multiple independent LoRA adapters concurrently on shared GPU hardware. These frameworks optimize GPU memory allocation, employing advanced paging mechanisms (e.g., PagedAttention) and custom CUDA kernels (e.g., segment-gather GEMM) to route batches of requests to their respective adapters.

However, while these serving frameworks resolve the low-level GPU memory tiling and batch scheduling bottlenecks, they are completely agnostic to the temporal dynamics of recurrent dynamic merging. Specifically, they do not address how stateful routing states are decoupled or tracked across interleaved user requests. Our proposed Tenant-Decoupled Stateful Routing (TDSR) is directly orthogonal and complementary to these infrastructures: TDSR operates at the routing scheduler layer, maintaining and isolating small recurrent state slots to enable robust stateful model merging under multi-tenant workloads.
